% for option: --include-in-header
%
% Caution: the content will override metadata field 'header-includes'
% https://github.com/jgm/pandoc/issues/3139
% https://github.com/jgm/pandoc/issues/3138

% The automatically generated 'header-includes' metadata when Pandoc 
% does not use the '--include-in-header' option.
\makeatletter
\@ifpackageloaded{subfig}{}{\usepackage{subfig}}
\@ifpackageloaded{caption}{}{\usepackage{caption}}
\captionsetup[subfloat]{margin=0.5em}
\AtBeginDocument{%
\renewcommand*\figurename{Figure}
\renewcommand*\tablename{Table}
}
\AtBeginDocument{%
\renewcommand*\listfigurename{List of Figures}
\renewcommand*\listtablename{List of Tables}
}
\newcounter{pandoccrossref@subfigures@footnote@counter}
\newenvironment{pandoccrossrefsubfigures}{%
\setcounter{pandoccrossref@subfigures@footnote@counter}{0}
\begin{figure}\centering%
\gdef\global@pandoccrossref@subfigures@footnotes{}%
\DeclareRobustCommand{\footnote}[1]{\footnotemark%
\stepcounter{pandoccrossref@subfigures@footnote@counter}%
\ifx\global@pandoccrossref@subfigures@footnotes\empty%
\gdef\global@pandoccrossref@subfigures@footnotes{{##1}}%
\else%
\g@addto@macro\global@pandoccrossref@subfigures@footnotes{, {##1}}%
\fi}}%
{\end{figure}%
\addtocounter{footnote}{-\value{pandoccrossref@subfigures@footnote@counter}}
\@for\f:=\global@pandoccrossref@subfigures@footnotes\do{\stepcounter{footnote}\footnotetext{\f}}%
\gdef\global@pandoccrossref@subfigures@footnotes{}}
\newcommand*\listoflistings\lstlistoflistings
\AtBeginDocument{%
\renewcommand*{\lstlistlistingname}{List of Listings}
}
\makeatother

% Write code below

\usepackage{tikz}
\usetikzlibrary{
  arrows.meta,
  calc,
  math,
  intersections,
  through,
  angles,
  quotes,
  decorations,
  decorations.markings,
  cartes,
  positioning,
  3d,
}

\usepackage{xpatch}

\BeforeBeginEnvironment{tikzpicture}{\begin{center}}
\AfterEndEnvironment{tikzpicture}{\end{center}}

\lstset{
  % 使 `listings` 宏包遇到 `TeX` 时按照 `[LaTeX]TeX` 的规则来处理
  defaultdialect=[LaTeX]TeX,
}

\usepackage{bm}
 
\lstdefinestyle{mystyle}{
  frame=shadowbox,      % 设置边框样式，例如带阴影的边框
  frameround=tttt,      % 设置四个角都为圆角
  rulesepcolor=\color{gray}, % 设置阴影颜色
  backgroundcolor=\color{cyan},   
  commentstyle=\color{green},
  keywordstyle=\color{magenta},
  numberstyle=\tiny\color{gray},
  stringstyle=\color{purple},
  basicstyle=\footnotesize,
  breakatwhitespace=false,         
  breaklines=true,                 
  captionpos=b,                    
  keepspaces=true,                 
  numbers=left,                    
  numbersep=10pt,                  
  showspaces=false,                
  showstringspaces=false,
  showtabs=false,                  
  tabsize=2
}

\lstset{style=mystyle}